\section{Interpolating the width-uniform cold bound}
\label{app:spectral-interpolation}

The spectral residual argument gives a finite-scale theorem with a
threshold that does not depend on learner width. It can also sharpen
the dependence on the exponent $\gamma$ before the regime $\gamma\ge2$.
We record the interpolation explicitly, including its common sample
event.

Let $0<u_c\le1$ be a width-independent cold threshold, and let
$\mathcal G_c$ be its finite comparison class, with
$\log|\mathcal G_c|\le J_c$. Suppose the cold construction has the
following consequences for every $N\ge C$ and $0<t^2\le u_c$:
there is posterior mass at least $1-\beta$ whose members have
representatives in $\mathcal G_c$ with zero empirical ramp loss and
population ramp loss at least their test-margin error. In particular,
a law $\Gamma_c$ on $\mathcal G_c$ can be chosen with
\begin{equation}\label{eq:spectral-cold-interface}
 \widehat R_{\psi,S}(\Gamma_c)=0,\qquad
 R(Q_{S,N,t,c})\le\beta+R_\psi(\Gamma_c).
\end{equation}
The original posterior still satisfies
$Q_{S,N,t,c}\le\exp(NK/t^2)\mu_{N,t}$, where $K\ge1$ is any
valid target-cost bound, for example $D_*+1$.
Theorem~\ref{thm:spectral-cold} supplies this interface with the
threshold in \eqref{eq:spectral-cutoff}, the code budget in
\eqref{eq:spectral-code-budget}, and $\beta=\varepsilon/2$. The raw
target cost $K$ here is distinct from its residual-energy cutoff.

\begin{corollary}[Exponent interpolation from a uniform cold threshold]
\label{cor:spectral-interpolation}
Assume \eqref{eq:spectral-cold-interface}, and fix $\gamma>1$.
Put
\[
 B_\gamma=K u_c^{-\gamma/(\gamma-1)}.
\]
Let $M_\gamma,R_{\tau,\gamma}$ be the finite Gaussian reference-family
sizes of Lemma~\ref{lem:si-grid}, with its parameter $J=B_\gamma$, and
set
\[
 \Lambda_\gamma=1+
     \log\frac{2M_\gamma R_{\tau,\gamma}}{\delta}.
\]
With sample probability at least $1-\delta$, simultaneously for every
$N\ge C$ and $0<t^2\le N^{1-\gamma}$,
\begin{equation}\label{eq:spectral-interpolation-risk}
 R(Q_{S,N,t,c})\le\beta+
             \frac2n[J_c+B_\gamma+\Lambda_\gamma].
\end{equation}
Consequently, if $\beta\le\varepsilon/2$, it suffices to take
\begin{equation}\label{eq:spectral-interpolation-samples}
 n\ge\frac4\varepsilon
       [J_c+B_\gamma+\Lambda_\gamma].
\end{equation}
One may replace the sum $J_c+B_\gamma$ by their maximum after
adjusting the harmless finite-family confidence charge.
\end{corollary}
\begin{proof}
Use one reference prior with two equally weighted parts: the uniform
law on $\mathcal G_c$, and the finite Gaussian family from
Lemma~\ref{lem:si-grid}. Use the ramp loss on the first component and
the test-margin indicator on the Gaussian components. Apply the single
relative-entropy risk event of Lemma~\ref{lem:si-risk}.

Write $u=t^2$. If $u\le u_c$, the law in
\eqref{eq:spectral-cold-interface} has zero empirical loss and relative
entropy at most $J_c+\log2$ under this reference prior. Its population
loss is at most $2[J_c+\log(2/\delta)]/n$.

If $u>u_c$ and $u\le N^{1-\gamma}$, then
\[
 N\le u^{-1/(\gamma-1)},\qquad
 NK/u\le K u^{-\gamma/(\gamma-1)}<B_\gamma.
\]
The Gaussian reference family therefore applies to the actual posterior.
Its empirical margin-indicator loss is zero, and its relative-entropy
charge is at most
$B_\gamma+1+\log(2M_\gamma R_{\tau,\gamma})$.
The resulting risk bound is at most
$2[B_\gamma+\Lambda_\gamma]/n$.
These two cases occur on the same sample event and prove
\eqref{eq:spectral-interpolation-risk}. The sample condition follows
by allocating at most $\varepsilon/2$ to the posterior tail and at
most $\varepsilon/2$ to the statistical term.
\end{proof}

For the Boolean target-cylinder budgets and the spectral residual
construction, at fixed $d,c,\varepsilon,\delta$,
\[
 K=\widetilde O(CL),\qquad
 u_c^{-1}=\widetilde O(CL^3),\qquad
 J_c=\widetilde O(C^4L^7).
\]
Thus for every fixed $\gamma>1$, the interpolated sufficient order is
\begin{equation}\label{eq:spectral-interpolation-powers}
 \widetilde O\left(
 C^4L^7+
 C^{\,2+1/(\gamma-1)}L^{\,4+3/(\gamma-1)}\right).
\end{equation}
Taking the better of this and the all-scale weighted-code theorem gives
\[
 \widetilde O\!\left(\min\left\{
 C^4L^9,\;
 C^4L^7+C^{\,2+1/(\gamma-1)}L^{\,4+3/(\gamma-1)}
 \right\}\right).
\]
The choice of which sufficient bound to use can be made from the resource
budgets before the sample is drawn. No sample-dependent choice of a
reference prior is required.

In particular, $\gamma\ge7/4$ permits
$\widetilde O(C^4L^8)$, and $\gamma\ge2$ permits
$\widetilde O(C^4L^7)$. These statements can use constants uniform in
the indicated exponent ranges. At the endpoint $\gamma=7/4$, the
more precise two-term order is
$\widetilde O(C^4L^7+C^{10/3}L^8)$.
The separate multilevel bound
$\widetilde O(\min\{C^4L^7,C^{9/2}L^6\})$ remains available in its
stated stronger-scale regimes.

For explicit logarithms, write $F_0=1+\log^2(2+C)$ and
$G=\log(2+CL)$. The target and cold-threshold estimates give
\[
 B_\gamma=
 O\!\left(C^{\,2+1/(\gamma-1)}
          L^{\,4+3/(\gamma-1)}F_0^{\,2+1/(\gamma-1)}\right),
\]
with fixed $\gamma>1$ and fixed ancillary resources.
For $\gamma\ge7/4$, a common log-retaining upper bound is
\[
 O\!\left(C^4L^8F_0^4\log^2(2+L)G\right).
\]
The finite-family charge $\Lambda_\gamma$ is logarithmic in the
resource budgets; it is uniformly $O(G)$ when $\gamma\ge7/4$ at the
fixed confidence parameter. The all-scale theorem has its own
$F_0^4G^2$ factor and remains preferable when the displayed interpolation
is larger.
